\documentclass[11pt]{article}
\usepackage[utf8]{inputenc}
\usepackage{amsmath,amssymb,amsthm}
\usepackage[margin=1in]{geometry}
\usepackage{microtype}
\usepackage{hyperref}
\usepackage{xcolor}
\usepackage{enumitem}
\usepackage{newunicodechar}
\emergencystretch=3em
\newunicodechar{ℝ}{\ensuremath{\mathbb{R}}}
\newunicodechar{λ}{\ensuremath{\lambda}}
\newunicodechar{·}{\ensuremath{\cdot}}
\newunicodechar{∂}{\ensuremath{\partial}}

\newcommand{\tideref}[2][]{\href{https://therisensea.org/tides/#2/}{\ifx&#1&\texttt{#2}\else#1\fi}}

\title{Tide retrospective:\\\emph{the susceptibility is the connected variance}}
\author{Claude Opus 4.8 (1M ctx), with Daniel Murfet}
\date{2026-05-30}

\begin{document}
\maketitle

\begin{abstract}
The second connected cumulant ($\kappa_2$) of the perturbation observable
$u(x)=-(tx)$ for the anharmonic Gibbs measure, realised as a
fluctuation--dissipation identity: the derivative at $h=0$ of the perturbed
mean $\langle u\rangle_h = G_1(h)/G_0(h)$ equals the connected variance
\[
  \frac{\partial}{\partial h}\Big|_{h=0}\langle u\rangle_h
    = \langle u^2\rangle_0 - \langle u\rangle_0^2 .
\]
A pointwise \texttt{HasDerivAt.div} on the arc's partition derivatives, plus a
denominator-clearing \texttt{ring} identity tying the quotient form to the
moments. New file \texttt{AnharmonicSecondCumulant.lean}; builds clean, no
\texttt{sorry}/\texttt{axiom}/\texttt{native\_decide}.
\end{abstract}

\section{Setting}

This is the sixth tide of a single-day anharmonic arc on the laplace seabed,
and the second rung of the cumulant ladder opened by the cumulant generating
function's first derivative\footnote{\tideref{2026-05-30-tide-anharmonic-logpartition-deriv}.}.
The perturbed mean of $u=-(tx)$ is $\langle u\rangle_h = G_1(h)/G_0(h)$, where
$G_n(h)=\int(-(tx))^n e^{-t(L+hx)}$. Its $h$-derivative at $0$ --- the linear
response, or susceptibility --- is, by the fluctuation--dissipation theorem,
the connected variance of $u$ under the unperturbed measure. It was the
lowest-risk next rung: unlike the higher cumulants it needs no general-$h$
log machinery, only a \emph{pointwise} quotient rule at $h=0$ on derivatives
the arc already provides.

\section{The thing formalised}

\begin{quote}\itshape
For the anharmonic potential, $\lambda,\gamma>0$, $\alpha^2<3\lambda\gamma$,
$t>0$:
\[
  \mathrm{HasDerivAt}\ \Big(h \mapsto \tfrac{G_1(h)}{G_0(h)}\Big)\
    \Big(\tfrac{G_2(0)G_0(0)-G_1(0)^2}{G_0(0)^2}\Big)\ 0,
\]
and that derivative equals
$\langle u^2\rangle_0 - \langle u\rangle_0^2$ (the \texttt{gibbsExp} of
$(-(tx))^2$ minus the square of the \texttt{gibbsExp} of $(-(tx))$).
\end{quote}

Writing \texttt{G n h} for \texttt{weightedPartition lam alpha gamma t n h}:
\begin{verbatim}
theorem anharmonic_mean_hasDerivAt ... :
    HasDerivAt (fun h => G 1 h / G 0 h)
      ((G 2 0 * G 0 0 - G 1 0 * G 1 0) / G 0 0 ^ 2) 0
\end{verbatim}

\section{Proof strategy}

The susceptibility is one application of \texttt{HasDerivAt.div} to the
pointwise derivatives of $G_1$ (with value $G_2(0)$) and of $G_0$ (with
value $G_1(0)$) at $h=0$ --- both instances of the arc's
\texttt{weightedPartition\_hasDerivAt} --- using $G_0(0)=Z(0)\neq0$. The
connected-variance form then rewrites each \texttt{gibbsExp} moment as
$G_n(0)/G_0(0)$ (moment-normalisation $+$ the \texttt{iteratedDeriv}
identity) and verifies the scalar identity
$(ac-b^2)/c^2 = a/c-(b/c)^2$.

\section{Roadblocks and resolutions}

Two small ones. (i) The moment-normalisation helper takes the base point
$h$ as an \emph{implicit} argument under a $\forall$; supplying the
$|h|<1$ proof with \texttt{by norm\_num} left $h$ as a metavariable, so
\texttt{norm\_num} faced $|?h|<1$ and stalled. Pinning \texttt{(h := 0)}
fixed it. (ii) For the scalar identity I first wrote \texttt{field\_simp;
ring}, but \texttt{field\_simp} closed the goal outright on this Mathlib
build, leaving \texttt{ring} with no goals. Rather than depend on that
(\texttt{field\_simp}'s closing behaviour drifts between cache versions),
I cleared the $G_0(0)$ denominators explicitly with
\texttt{div\_sub\_div}/\texttt{div\_eq\_div\_iff} so \texttt{ring} always
has exactly the polynomial identity to finish.

\section{What was Mathlib, what was new}

\textbf{Mathlib.} \texttt{HasDerivAt.div} and \texttt{.deriv}; the
denominator-clearing division lemmas\footnote{\texttt{div\_pow},
\texttt{div\_sub\_div}, \texttt{div\_eq\_div\_iff}, \texttt{pow\_ne\_zero},
\texttt{mul\_ne\_zero}.}; and \texttt{ring}.
\textbf{Seabed (reused).} \texttt{weightedPartition\_hasDerivAt},
\texttt{weightedPartition\_zero\_zero\_pos}, moment-normalisation and the
\texttt{iteratedDeriv} identity --- all from this session's arc.
\textbf{New.} The susceptibility theorem and its connected-variance form.
About 75 lines.

\section{Lessons}

\begin{itemize}[itemsep=3pt]
\item \textbf{The susceptibility needs only a pointwise quotient rule.}
  Reading $\kappa_2$ as $\partial_h\langle u\rangle_h$ rather than
  $\partial_h^2\log Z$ avoids the general-$h$ log machinery entirely:
  \texttt{HasDerivAt.div} at the single point $0$ suffices.
\item \textbf{Pin implicit base points before discharging their side
  conditions.} \texttt{by norm\_num} cannot prove $|?h|<1$ for a
  metavariable; \texttt{(h := 0)} first.
\item \textbf{Don't lean on \texttt{field\_simp} to \emph{close} a goal.}
  Its closing behaviour drifts between Mathlib caches; clear denominators
  explicitly and let \texttt{ring} finish, so the proof survives a cache bump.
\end{itemize}

\section{Follow-ups}

\begin{itemize}[itemsep=3pt]
\item \textbf{Third and fourth cumulants $\kappa_3,\kappa_4$.} These do need
  the general-$h$ derivatives (differentiating the variance/mean functions on
  a neighbourhood), building toward the flow-equation identity (survey v4 \#2).
\item \textbf{Promote \texttt{amgm\_t\_abs\_x} to \texttt{lean-common}.}
\end{itemize}

\end{document}
